\documentclass[11pt,a4paper]{article}
\usepackage[T1]{fontenc}
\usepackage[utf8]{inputenc}
\usepackage{amsmath,amssymb,xcolor}
\usepackage[margin=25mm]{geometry}
\usepackage[colorlinks=true,urlcolor=blue]{hyperref}
\definecolor{nexorablue}{HTML}{2459B5}
\setlength{\parindent}{0pt}
\setlength{\parskip}{7pt}
\setlength{\emergencystretch}{2em}
\renewcommand{\familydefault}{\sfdefault}
\begin{document}

{\small\bfseries\color{nexorablue} NEXORA CPGE}\par
{\small MPSI / NEX-CIR-02 / Corrigé}\par
{\LARGE\bfseries Deux sources opposées}\par
\vspace{0.25em}\hrule\vspace{0.6em}

\section*{Équation de maille}
La somme des variations de potentiel est $18-4i-6-2i=0$. Il vient $i=2\,\mathrm A$ dans le sens choisi.

\section*{Bilan de puissance}
La source de $18\,\mathrm V$ reçoit $-36\,\mathrm W$ et fournit donc de l’énergie. La source de $6\,\mathrm V$ reçoit $12\,\mathrm W$. Les deux résistances dissipent $16$ et $8\,\mathrm W$. On vérifie $-36+12+16+8=0$.


\end{document}
